\section*{Using This Manual}
Start with Quick Start to get a signal on screen. The numbered panel drawing
is a map to the control reference. The walkthroughs then turn those controls
into experiments: patch a source, choose a known setup, change one thing,
and interpret the result.

\begin{center}
\renewcommand{\arraystretch}{1.3}
\begin{tabularx}{\textwidth}{@{}Xr@{}}
\toprule
\textbf{Find your way} & \textbf{Page} \\
\midrule
\hyperref[sec:overview]{Overview} & \pageref{sec:overview} \\
\hyperref[sec:quick-start]{Quick Start} & \pageref{sec:quick-start} \\
\hyperref[sec:panel]{Panel Layout And Controls} & \pageref{sec:panel} \\
\hyperref[sec:display]{Reading The Display} & \pageref{sec:display} \\
\hyperref[sec:settings]{Context Menu And Saved Settings} & \pageref{sec:settings} \\
\hyperref[sec:shapes]{Recognizing Sound Shapes} & \pageref{sec:shapes} \\
\hyperref[sec:envelope]{Following An Envelope} & \pageref{sec:envelope} \\
\hyperref[sec:detail]{Finding Quiet Detail} & \pageref{sec:detail} \\
\hyperref[sec:time]{Time And Measurement Limits} & \pageref{sec:time} \\
\hyperref[sec:presets]{Factory Presets} & \pageref{sec:presets} \\
\hyperref[sec:troubleshooting]{Troubleshooting} & \pageref{sec:troubleshooting} \\
\hyperref[sec:further-reading]{Analysis Timing And Further Reading} & \pageref{sec:further-reading} \\
\bottomrule
\end{tabularx}
\end{center}

\paragraph{Try the experiments}
Use any Rack modules that provide the named functions; no particular
oscillator, filter, or envelope is required. Settings in the walkthroughs
are starting points, not new factory presets. Values in milliseconds and
hertz assume the stated Rack sample rate. Keep your listening path connected
separately: the analyzer has no audio outputs.

\paragraph{Keep comparisons fair}
Use the same source and sample rate, match input gains, and set Slope to zero
before comparing levels. Change one control at a time. Schematic figures
explain relationships rather than reproduce captured data; the cover shows
the rendered module.
